% !TeX program = xelatex
% UTF-8. Standalone: run xelatex CAP.tex from this directory.
% For inclusion in a XePersian document, define \CAPInNotebook first.
% Source: Google Drive Notebook/CAP, seven images in filename order.
% pic01.jpg = 20261010_052026465.jpg; 1grnZAk674ew8WSr0ZLsXN8Myy70NKRfx
% pic02.jpg = 20261010_052033038.jpg; 1wq4Q4ihQKSq9Wma-GOSKdGO9cou-zTA4
% pic03.jpg = 20261010_052037483.jpg; 13Q6sYJsGhy3mtc34dvdYOPfyZxTnqBju
% pic04.jpg = 20261010_052043927.jpg; 1JDmx8czJC2H3yvAZeHAwKpAVP-CWoIbr
% pic05.jpg = 20261010_052050440.jpg; 1bBxeIlFMaSuaD2TOKWD75fKDjDWgBhxQ
% pic06.jpg = 20261010_052057841.jpg; 1be2-HIZ52rH9XPVELzLgUQaWr0AicFVy
% pic07.jpg = 20261010_052102313.jpg; 1ONJC9zNa4fg168Ov5jDjVKrG2jGKJHKP
\ifdefined\CAPInNotebook
  \let\CAPfinish\relax
\else
  \def\CAPfinish{\end{document}}
  \documentclass[12pt,a4paper]{article}
  \usepackage[margin=22mm]{geometry}
  \usepackage{amsmath,amssymb}
  \usepackage{graphicx}
  \usepackage{array,longtable}
  \usepackage{xepersian}
  \IfFontExistsTF{Amiri}{\settextfont{Amiri}}{\settextfont{DejaVu Sans}}
  \setlatintextfont{DejaVu Serif}
  \setlength{\parindent}{0pt}
  \setlength{\parskip}{5pt}
  \setlength{\emergencystretch}{3em}
  \begin{document}
\fi
\providecommand{\CAPunclear}[1]{\textbf{[خوانش نامطمئن: #1]}}
\providecommand{\CAPpage}[2]{\clearpage\section*{تصویر #1}\lr{\texttt{#2}}\par}
\providecommand{\CAPmath}[1]{\lr{\(#1\)}}

\begin{center}
{\Large\bfseries \lr{CAP}}\par
\lr{Nelson 22nd}\par
۴۴۹
\end{center}
\textbf{یادداشت رونویسی:} متن به ترتیب هفت تصویر، با حفظ عبارت‌ها، اختصارات و اعداد منبع نوشته شده است. فلش‌ها به صورت رابطه یا فهرست بازنویسی شده‌اند. خوانش‌های نامطمئن صریحاً مشخص شده‌اند؛ اعداد یا گزاره‌های پزشکی منبع اصلاح نشده‌اند. تصاویر کامل در پیوست آمده‌اند.

\CAPpage{۱}{pic01.jpg}
\textbf{سربرگ و حاشیه:} \lr{CAP}، ۴۴۹؛ \lr{Subject: Nelson 22nd}؛ \lr{Date: / /}؛ شمارهٔ صفحه: ۱؛ نشان پایین برگه: \lr{YASHA}.

پنومونی خصوصاً در کودکان زیر ۵ سال می‌تواند کشنده باشد.

مرگ ناشی از پنومونی با فقر مرتبط است.
\subsection*{اتیولوژی پنومونی}
\begin{itemize}
\item میکروب‌ها.
\item آسپیراسیون.
\item واکنش‌های دارویی یا حساسیت.
\item پنومونیت ناشی از تشعشع.
\item پنوموکوک شایع‌ترین باکتری عامل پنومونی در ۳ هفته تا ۵ سال است.
\item مایکوپلاسما پنومونیه و کلامیدوفیلا پنومونیه شایع‌ترین پاتوژن باکتری‌ها در ۵ سال و بیشتر.
\item سایر باکتری‌ها: \lr{GAS} (استرپ پایوژن)، استاف اورئوس.
\end{itemize}
به طور شایع استاف اورئوس و استرپ پنومونیه، روی آنفولانزا سوار می‌شوند.
\begin{itemize}
\item در کودک مبتلا به \lr{HIV}، علاوه بر موارد بالا: \lr{TB}، مایکوباکتریا غیر \lr{TB}، سالمونلا، \lr{E. Coli}، پنوموسیستیس جیرووسی، \lr{CMV}.
\item شایع‌ترین علل کلی پنومونی بالای ۱ ماه تا ۵ سال: ویروس‌ها.
\item در ۳۰\% موارد، پنومونی‌های ویروسی، چند ویروس دارند.
\item \lr{RSV} و رینوویروس شایع‌ترین ویروس‌ها؛ خصوصاً زیر ۲ سال.
\item پنومونی‌های ویروسی در پاییز و زمستان بسیار بسیار شایع‌تر هستند.
\end{itemize}

\CAPpage{۲}{pic02.jpg}
\textbf{سربرگ و حاشیه:} ۴۴۹؛ \lr{Subject:} (خالی)؛ \lr{Date: / /}؛ شمارهٔ صفحه: ۲؛ \lr{YASHA}.

پنومونی با سودومونا در کودک با \lr{CF} شایع‌تر است.
\subsection*{تظاهرات بالینی در پنومونی}
\begin{itemize}
\item معمولاً از چند روز جلوتر علائم \lr{URI} وجود دارد؛ خصوصاً رینیت و سرفه.
\item در پنومونی وایرال تب وجود دارد ولی معمولاً پایین‌تر از باکتریال است.
\item مهم‌ترین یافته تاکی‌پنه است.
\item افزایش کار تنفسی: رترکشن بین‌دنده‌ای، زیر دنده‌ای، سوپرااسترنال، نازال فلرینگ، استفاده از عضلات فرعی.
\item در موارد شدید: سیانوز و لتارژی؛ خصوصاً در شیرخواران.
\item سمع ریه: کراکل و ویزینگ.
\item معمولاً افتراق باکتریال از وایرال با معاینه ممکن نیست.
\item پنومونی باکتریال در کودکان بزرگ‌تر معمولاً به صورت شروع ناگهانی تب بالا، سرفه و درد قفسهٔ سینه است.
\item دلیریوم، اضطراب ممکن است.
\item کودکان معمولاً به سمت ریهٔ مبتلا می‌خوابند تا درد پلورتیک کم شود.
\item پنومونی لوب تحتانی ممکن است درد شکم ایجاد کند ولی تندرنس ندارد.
\item دیستنشن شکم به دلیل پرهوا شدن معده به دلیل بلع هوا.
\item ایلئوس در اواخر.
\item ممکن است کبد به صورت اشتباه بزرگ لمس شود، چون دیافراگم به دلیل پرهوایی ریه به آن فشار می‌آورد.
\end{itemize}
سمع خصوصاً در کودکان کم‌سن ممکن است گمراه‌کننده باشد.

تب نسبت به تاکی‌پنه کمتر مهم است.

\CAPpage{۳}{pic03.jpg}
\textbf{سربرگ و حاشیه:} ۴۴۹؛ \lr{Subject:} (خالی)؛ \lr{Date: / /}؛ شمارهٔ صفحه: ۳؛ \lr{YASHA}.

برخی شیرخواران با پنومونی باکتریال: علائم درگیری \lr{GI}:
\begin{itemize}
\item تهوع، بی‌اشتهایی، اسهال، دیسترس شکمی، ایلئوس پارالیتیک.
\end{itemize}
سیانوز مطرح‌کنندهٔ درگیری مولتی‌لوبار است.
\subsection*{ریسک فاکتورهای پنومونی \lr{Severe}}
\begin{itemize}
\item \CAPmath{T>38.5}.
\item تاکی‌پنه.
\item رترکشن.
\item نازال فلرینگ.
\item گرانتینگ.
\item کپیلاری ریفیل: \CAPmath{2<} ثانیه. \CAPunclear{علامت در تصویر بعد از عدد و رو به چپ نوشته شده؛ جهت مقایسه عین جایگاه علامت حفظ شده است.}
\item سیانوز.
\item تاکی‌کاردی.
\item \lr{poor feeding}.
\end{itemize}
\subsection*{تشخیص پنومونی}
\begin{itemize}
\item مشاهدهٔ اینفیلتریشن در \lr{CXR} (۲ نما) از تشخیص پنومونی حمایت می‌کند.
\item پنومونی وایرال، در \lr{CXR}: شواهد پرهوایی، اینفیلتراسیون دوطرفهٔ اینترستیشیال، پری‌برونکیال \lr{cuffing}.
\item درگیری واضح یک لوب در \lr{CXR} معمولاً در پنوموکوک دیده می‌شود.
\item \lr{CXR} نمی‌تواند اتیولوژی پنومونی را تشخیص دهد.
\item در پنومونی بدون عارضه تکرار \lr{CXR} برای تأیید درمان شدن نیاز نیست.
\item در بیمار مشکوک به پنومونی (تاکی‌پنه، سرفه، تب، کراکل دوطرفه در کاهش صدا تنفسی) که تحت درمان سرپایی درمان کنیم نباید \lr{CXR} بگیریم. \CAPunclear{حرف ربط پیش از «کاهش صدا تنفسی» و عبارت «تحت درمان سرپایی درمان کنیم».}
\end{itemize}
\lr{POCUS} بسیار حساس است برای تشخیص کانسالیدیشن، ایر برونکوگرام و ایفیوژن.

\CAPpage{۴}{pic04.jpg}
\textbf{سربرگ و حاشیه:} ۴۴۹؛ \lr{Subject:} (خالی)؛ \lr{Date: / /}؛ شمارهٔ صفحه: ۴؛ \lr{YASHA}.
\begin{itemize}
\item در پنومونی وایرال \lr{WBC} ممکن است \lr{N/L} یا افزایش باشد ولی معمولاً بالاتر از ۲۰ نیست و لنف دامیننت است.
\item در پنومونی باکتریال \lr{WBC} معمولاً بین \lr{15--40} و \lr{PMN} دامیننت است. \CAPunclear{حد بالای بازه «۴۰» خوانده می‌شود.}
\end{itemize}
تشخیص قطعی پنومونی وایرال: شناسایی ویروس، \lr{PCR} است.

تشخیص قطعی پنومونی باکتریال: ایزوله کردن باکتری در خون، مایع پلور یا ریه.
\begin{itemize}
\item کشت خلط ارزش کمی دارد در کودکان.
\item کشت خون تنها در ۱۰\% پنومونی‌های باکتریال مثبت می‌شود.
\item استرپ \lr{GAS} و هموفیلوس باکتریمی می‌دهند.
\item کشت در بیمار سرپایی و \lr{non-toxic} انجام ندهیم.
\item سیر بهبودی ندارد، دتریوریت می‌شود، عارضه‌دار شده: نیاز به بستری؛ \lr{B/C}. \CAPunclear{واژهٔ انگلیسیِ فارسی‌نویسی‌شده «دتریوریت».}
\end{itemize}
\subsection*{تشخیص پنومونی}
\begin{itemize}
\item \lr{PCR} یا کشت نمونهٔ نازوفارنکس.
\item برای تشخیص \lr{GAS} سرولوژی: تست آنتی‌استرپتولایزین \lr{O}، آنتی \lr{DNase B}.
\end{itemize}
\subsection*{فاکتورهای مطرح‌کنندهٔ نیاز به بستری}
\begin{itemize}
\item سن زیر ۶ ماه.
\item ایمنی مختل.
\item ظاهر \lr{toxic}.
\item دیسترس تنفسی متوسط تا شدید.
\item سیانوز و هایپوکسی \CAPmath{(<90\%)}.
\item شوک.
\item پنومونی عارضه‌دار.
\item تشنج، عدم تحمل \lr{PO}.
\item سندرم سیکل سل، سندرم \lr{Acute chest}.
\item دهیدریشن شدید.
\item عدم پاسخ به درمان مناسب \lr{AB}.
\item عوامل اجتماعی.
\item پاتوژن‌های \lr{high risk}.
\end{itemize}

\CAPpage{۵}{pic05.jpg}
\textbf{سربرگ و حاشیه:} ۴۴۹؛ \lr{Subject:} (خالی)؛ \lr{Date: / /}؛ شمارهٔ صفحه: ۵؛ \lr{YASHA}.
\subsection*{درمان پنومونی}
\begin{itemize}
\item کودک \lr{mildly ill}: درمان سرپایی؛ آموکسی‌سیلین:
\begin{latin}
90 mg/kg/day orally divided twice daily.
\end{latin}
\item آلترناتیو: سفوروکسیم، آموکسی/کلاو.
\item کودک در سن مدرسه و نوجوانان (مایکوپلاسما و \lr{C.}): ماکرولید؛ آزیترومایسین.
\item آلترناتیو: کلاریترومایسین.
\item داکسی‌سایکلین: \CAPmath{8<} سال. \CAPunclear{علامت مقایسه در تصویر پس از عدد ۸ است؛ جایگاه آن حفظ شده است.}
\item نوجوانان: فلوروکینولون تنفسی (لوو، موکسی).
\item آلترناتیو ماکرولید هستند.
\end{itemize}
درمان تجربی کودک بستری باید بر اساس اپیدمیولوژی لوکال باشد.
\begin{itemize}
\item در منطقه‌ای که مقاومت بالا به پنی‌سیلین در استرپ پنومونیه وجود ندارد، کودکی که کاملاً علیه هموفیلوس تیپ \lr{B} و استرپ پنومونیه واکسینه شده است، در \lr{Severely ill} نمی‌باشد: آمپی‌سیلین یا پنی‌سیلین \lr{G}.
\item سایر کودکان: سفتریا، سفوتاکسیم.
\item اگر شک به مایکوپلاسما، \lr{C.} و غیره در موارد بالا: ماکرولید هم \lr{add} می‌کنیم.
\item اگر شک به استافیلوکوک (پنوماتوسل، امپیم): ونکو یا کلیندا \lr{added} می‌شود.
\item کودک با نارسایی تنفسی در زمینهٔ آنفولانزا که \lr{MRSA} روی آن سوار شده: ونکو + کلیندا.
\end{itemize}

\CAPpage{۶}{pic06.jpg}
\textbf{سربرگ و حاشیه:} ۴۴۹؛ \lr{Subject:} (خالی)؛ \lr{Date: / /}؛ شمارهٔ پایین تصویر ظاهراً ۹ (با وجود ششم بودن تصویر)؛ \lr{YASHA}.
\begin{itemize}
\item اگر شک به پنومونی وایرال داریم و کودک \lr{mildly ill} است آنتی‌بیوتیک نمی‌دهیم (کودک نباید در دیسترس باشد).
\item اگر سیر بیماری بهتر نشد: شروع آنتی‌بیوتیک.
\end{itemize}
\subsection*{طول درمان آنتی‌بیوتیکی}
به آنقدری باشد که کودک ۷۲ ساعت بدون تب باشد.
\begin{itemize}
\item طول دورهٔ ۵–۷ روزه خصوصاً برای موارد سرپایی مناسب است.
\item درمان طولانی‌مدت در پنومونی بدون عارضه سودی ندارد.
\item برخی مطالعات گفته‌اند تا رسیدن پروکلسی‌تونین به \lr{0.1--0.25} میکروگرم در لیتر.
\end{itemize}
\subsection*{درمان خوراکی با \lr{Zinc}}
\begin{itemize}
\item سن زیر ۱۲ ماه: \lr{10 mg} روزانه تا ۷ روز.
\item سن ۱۲ ماه و بالاتر: \lr{20 mg} روزانه تا ۷ روز.
\item در کشورهای با درآمد متوسط و ضعیف، در پنومونی تا \lr{Severe}. \CAPunclear{عبارت پایانی «تا Severe».}
\end{itemize}
\subsection*{پروگنوز پنومونی}
ظرف ۴۸–۷۲ ساعت علائم شروع به بهبودی می‌کنند: تب، سرفه، تاکی‌پنه.

زمانی که بیمار علی‌رغم آنتی‌بیوتیک مناسب بهتر نمی‌شود:
\begin{enumerate}
\item عارضه: ایفیوژن، امپیم.
\item مقاومت آنتی‌بیوتیکی.
\item اتیولوژی باکتری نیست.
\item انسداد برونشی به هر علتی.
\item بیماری‌های زمینه: نقص ایمنی، \lr{CF}.
\item علل غیرعفونی: پنومونیت (آسپیریشن، دارو، \ldots). \CAPunclear{واژهٔ نخست داخل پرانتز «آسپیریشن» خوانده می‌شود.}
\end{enumerate}

\CAPpage{۷}{pic07.jpg}
\textbf{سربرگ و حاشیه:} ۴۴۹؛ \lr{Subject:} (خالی)؛ \lr{Date: / /}؛ شمارهٔ صفحه: ۷؛ \lr{YASHA}.
\begin{itemize}
\item در مواجهه با عدم پاسخ، قدم اول: \lr{CXR}.
\item در موارد نارسایی تنفسی ممکن است \lr{BAL} اندیکاسیون داشته باشد.
\item \lr{HRCT} ممکن است نیاز شود.
\end{itemize}
\subsection*{عارضه‌های پنومونی}
\begin{itemize}
\item ایفیوژن پاراپنومونیک: استاف اورئوس، استرپ پنومونیه، استاف پایوژن (\lr{GAS}).
\item اغلب ایفیوژن‌ها استریل هستند.
\item در لترال دکوبیتوس: پلورال ایفیوژن پاراپنومونیک اگر کمتر از \lr{1 cm} باشد نیاز به تخلیه نیست.
\end{itemize}
\subsection*{افتراق \lr{PE} اگزوداتیو از ترانسوداتیو}
شمارهٔ دست‌نویس کنار جدول: ۱).

\textbf{یادداشت جدول:} عنوان ستون‌ها عین تصویر \lr{T} و \lr{E} است. جایگاه و جهت علامت‌های مقایسه، حتی وقتی پس از عدد نوشته شده‌اند، حفظ شده است. خانهٔ آخر ستون \lr{E} در منبع خالی است.
\begin{latin}
\renewcommand{\arraystretch}{1.35}
\begin{longtable}{|p{0.40\linewidth}|p{0.24\linewidth}|p{0.24\linewidth}|}
\hline
Parameter & T & E \\
\hline
Glucose & \(40\leq\) & \(<40\) \\
\hline
\(LDH_{PE}/LDH_{plasma}\) & \(<0.6\) & \(0.6<\) \\
\hline
LDH & \(<200\) IU/day & \(200\) IU \(<\) \\
\hline
\(Pro_{PE}/Pro_{serum}\) & \(<0.5\) & \(0.5\leq\) \\
\hline
WBC & \(<10{,}000\) & \(500{,}000<\) \\
\hline
pro & \(<3\) gr/dL & \\
\hline
\end{longtable}
\end{latin}
\CAPunclear{واحد LDH در ستون T شبیه «IU/day» است؛ عدد WBC در ستون E به صورت «500,00» با یک صفر زیر آن نوشته شده و «500,000» خوانده شده است. این مقادیر اصلاح پزشکی نشده‌اند.}

پنومونی راجعه: حداقل ۲ اپیزود طی ۳ ماه؛ احتمال یک علت زمینه‌ای وجود دارد.
\CAPunclear{بازهٔ زمانی در دست‌نویس «۳ ماه» خوانده می‌شود.}

\clearpage
\section*{پیوست: هفت تصویر کامل منبع}
تمام حاشیه‌ها، سربرگ‌ها، اصلاحات و جدول اصلی در تصاویر زیر محفوظ‌اند.
% Compile from CAP/ or the repository root. Both paths use the corrected directory name.
\IfFileExists{../../ped-infectious diseases/CAP/images/pic01.jpg}{%
  \def\CAPimagebase{../../ped-infectious diseases/CAP/images/}}{%
  \def\CAPimagebase{pediatrics/ped-infectious diseases/CAP/images/}}
\newcommand{\CAPoriginal}[2]{%
  \clearpage\subsection*{اصل تصویر #1: \lr{\texttt{pic#2.jpg}}}%
  \begin{center}
  \includegraphics[width=\linewidth,height=0.87\textheight,keepaspectratio]{\CAPimagebase pic#2.jpg}
  \end{center}}
\CAPoriginal{۱}{01}
\CAPoriginal{۲}{02}
\CAPoriginal{۳}{03}
\CAPoriginal{۴}{04}
\CAPoriginal{۵}{05}
\CAPoriginal{۶}{06}
\CAPoriginal{۷}{07}
\CAPfinish
